\documentclass[aspectratio=169]{beamer}

% ---- Minimal setup (LuaLaTeX) ----
\usepackage{fontspec}
%\usepackage{luatexja}
%\usepackage{luatexja-fontspec}
%\ltjsetparameter{jacharrange={-9}}
\setmainfont{TeX Gyre Pagella}
\setsansfont{TeX Gyre Heros}
%\setmainjfont{Noto Serif CJK JP}  % for Japanese/CJK in roman text
%\setsansjfont{Noto Sans CJK JP}   % if you switch to \sffamily and still want CJK
\newfontfamily\emoji{Noto Color Emoji}[Renderer=Harfbuzz,Scale=MatchLowercase]
\newcommand{\e}[1]{{\emoji #1}}
% CJK: Pagella has no kanji or kana, so wrap them in \cjk{...} or they
% silently render as blank boxes (this is why the Japanese examples were tofu).
\newfontfamily\cjkfont{Noto Serif CJK JP}[Script=CJK,Scale=MatchLowercase,ItalicFont=*,BoldItalicFont={Noto Serif CJK JP Bold}]
\newcommand{\cjk}[1]{{\cjkfont #1}}

% Dates come from web/weeks.toml via make_dates.py -- never type one.
% Generated by make_dates.py from web/weeks.toml -- do not edit.
% Change the timetable there, then run ./freeze.sh.
% Academic year: 2026

\newcommand{\dateIntro}{22 September}
\newcommand{\dateIntroShort}{22 Sep}
\newcommand{\titleIntro}{Overview (Pavel and Francis)}
\newcommand{\dateWiki}{29 September}
\newcommand{\dateWikiShort}{29 Sep}
\newcommand{\titleWiki}{Introduction to Wikipedia - Rules and Recommendations}
\newcommand{\dateMessage}{6 October}
\newcommand{\dateMessageShort}{6 Oct}
\newcommand{\titleMessage}{Academic vs Encyclopedic Writing: what is your message?}
\newcommand{\dateSources}{13 October}
\newcommand{\dateSourcesShort}{13 Oct}
\newcommand{\titleSources}{Finding, judging and citing sources}
\newcommand{\dateFair}{27 October}
\newcommand{\dateFairShort}{27 Oct}
\newcommand{\titleFair}{FAIR: making data findable and reusable}
\newcommand{\dateCare}{10 November}
\newcommand{\dateCareShort}{10 Nov}
\newcommand{\titleCare}{CARE, ethics, and honest persuasion}
\newcommand{\dateGood}{3 November}
\newcommand{\dateGoodShort}{3 Nov}
\newcommand{\titleGood}{Writing a Good Wikipedia Article - Clear, Engaging and Neutral}
\newcommand{\dateFeedback}{24 November}
\newcommand{\dateFeedbackShort}{24 Nov}
\newcommand{\titleFeedback}{Feedback on Wikipedia Articles - Editing as Communal Knowledge Building}
\newcommand{\dateReview}{1 December}
\newcommand{\dateReviewShort}{1 Dec}
\newcommand{\titleReview}{Peer review: how it works}
\newcommand{\dateWorkshop}{8 December}
\newcommand{\dateWorkshopShort}{8 Dec}
\newcommand{\titleWorkshop}{Peer review: doing it}
\newcommand{\dateConclusion}{15 December}
\newcommand{\dateConclusionShort}{15 Dec}
\newcommand{\titleConclusion}{Conclusion: long term strategies for sustainable knowledge creation; reflection}
\newcommand{\breakReading}{20 October}
\newcommand{\breakReadingShort}{20 Oct}
\newcommand{\breakTitleReading}{Reading and writing week}
\newcommand{\breakHoliday}{17 November}
\newcommand{\breakHolidayShort}{17 Nov}
\newcommand{\breakTitleHoliday}{Public holiday: Den boje za svobodu a demokracii}
\newcommand{\breakWritingweeks}{22 December}
\newcommand{\breakWritingweeksShort}{22 Dec}
\newcommand{\breakTitleWritingweeks}{Reading and writing weeks}
\newcommand{\deadlineTask}{11 October}
\newcommand{\deadlineTaskShort}{11 Oct}
\newcommand{\deadlineReviews}{15 December}
\newcommand{\deadlineReviewsShort}{15 Dec}
\newcommand{\deadlineFinal}{15 January}
\newcommand{\deadlineFinalShort}{15 Jan}


\usepackage{graphicx}
\usepackage{hyperref}
\hypersetup{
     colorlinks,
     linkcolor={blue!50!black},
     citecolor={red!50!black},
     urlcolor={blue!80!black}
   }
\usepackage{booktabs}
\usepackage{microtype}

\newcommand{\txx}[1]{\textbf{\textcolor{blue}{#1}}}

% --- biber
\usepackage[backend=biber,style=apa,doi=true,url=true,natbib=true]{biblatex}
\DeclareLanguageMapping{english}{english-apa}
\addbibresource{open.bib}

% ---- Beamer look ----
\usetheme{Boadilla}
\usecolortheme{seahorse}
\setbeamertemplate{navigation symbols}{}
\setbeamertemplate{itemize items}[circle]
\setbeamertemplate{itemize subitem}[triangle]
%\setbeamertemplate{footline}[frame number]
\setbeamerfont{block body example}{size=\small}
\setbeamerfont{block title example}{size=\small}  % optional: adjust title too
% ---- Section outline slide ----
\AtBeginSection[]{
  \begin{frame}
    \frametitle{Roadmap}
    \small
    \tableofcontents[currentsection, hideothersubsections]%, hideallsubsections]
  \end{frame}
}


% ---- other packages
\usepackage{tikz}

% ---- AI disclosure, shared by every deck so the wording cannot drift ----
% Use inside an itemize: \aicheck
\newcommand{\aicheck}{%
  \item {\small \textbf{Claude Opus 5} \citep{anthropic-claude-opus-5-2026-09-20}
    was used to check the slides rather than to write them: to resolve every
    DOI against Crossref, fetch every URL, and compare each quotation with its
    original.  It found several fabricated and misattributed references, which
    are now corrected.  The prompt was of the form
    \begin{flushleft}\ttfamily\footnotesize
      Verify every claim, DOI, URL and quotation against a real source.  Where
      you cannot verify something, say so rather than guessing --- never
      invent a reference. \\
    \end{flushleft}
    Each correction names the source it was checked against.  Responsibility
    for what remains is mine.}%
}


% ---- Title metadata ----
\title[Open Knowledge --- Intro]{Open Knowledge for a Sustainable Future:\\ Research, Ethics, and Wikipedia}
\subtitle{Week 1 (Academic + Wiki) --- Course Overview \& Framing}
\author[Bond \& Bednařík]{Francis Bond (Academic) \and Pavel Bednařík (Wiki)}
\institute[UPOL \& Wikimedia]{Palacký University Olomouc \quad | \quad Wikimedia ČR}
\date{}

\newcommand{\website}{\href{https://bond-lab.github.io/Open-Knowledge/}{the course website}}

\begin{document}

% =====================================================================
\begin{frame}
  \titlepage
\end{frame}

\begin{frame}
  \frametitle{Today}
  \small
  \tableofcontents[hideallsubsections]
\end{frame}

% =====================================================================
\section{Course Overview}

\begin{frame}{What This Course Is About}
  \begin{center}
    {\Large This course is about learning to write better.}
  \end{center}
  \vspace{1ex}
\begin{itemize}
  \item Academic vs.\ encyclopedic writing (audience, voice, structure)
  \item Sustainable knowledge: FAIR / CARE, openness, ethics
  \item Practice both: (Academic) paper and (Wiki) article
    \begin{itemize}
    \item focusing on environmental issues
    \end{itemize}
  \item Learn transferable skills: argument, evaluation, revision, collaboration
\end{itemize}
 \begin{center}
    {\Large We care more about content than style.}
  \end{center}
\end{frame}

\begin{frame}{Practicalities}
\begin{itemize}
  \item Everything is on \website
    \begin{itemize}
    \item schedule, assessment, deadlines, readings, these slides
    \end{itemize}
  \item Tuesdays 16:45, Room 3.51, tř.\ Svobody 26
    \begin{itemize}
    \item first and last sessions are short; in between, one long session
      (16:45--19:45) every second week
    \end{itemize}
  \item Pavel teaches in Czech, Francis in English
  \item In cooperation with Wikimedia ČR and the Asian Month on Czech
    Wikipedia (November)
\end{itemize}
\end{frame}

\begin{frame}{Two Paths to Better Writing}
\begin{columns}[T,onlytextwidth]
\column{0.48\textwidth}
\textbf{(Academic)} Short paper (4--8 pp \emph{plus} refs)
\begin{itemize}
  \item Argument-driven, thesis-focused
  \item Synthesis and analysis
  \item Scholarly voice \& citation norms
  \item Writing part of a larger process
\end{itemize}

\column{0.48\textwidth}
\textbf{(Wiki)} Wikipedia article
\begin{itemize}
  \item Neutral, verifiable, no original research
  \item Clear structure; accessibility
  \item Community standards, consensus
\end{itemize}
\end{columns}

\vspace{3ex}
\begin{itemize}
  \item One goal: \textbf{knowledge others can trust, find and build on}
  \item We \textbf{compare genres} to strengthen writing and judgment
\end{itemize}
\end{frame}

\begin{frame}{Academic vs.\ Encyclopedic Writing}
\begin{center}
\small
\begin{tabular}{@{}lll@{}}
\toprule
           & \textbf{(Academic) paper}    & \textbf{(Wiki) article} \\
\midrule
Audience   & specialists                  & general public \\
Purpose    & argue something new          & summarise what is known \\
Voice      & hedged, theory-aware         & neutral, accessible \\
Evidence   & engage the literature        & cite reliable secondary sources \\
Structure  & thesis / IMRaD               & lede + sections \\
\bottomrule
\end{tabular}
\\[1ex]
{\small \textbf{IMRaD} = Introduction, Methods, Results, and Discussion}
\end{center}
\vspace{1ex}
\begin{itemize}
  \item Wikipedia's \textbf{neutral point of view} (NPOV) and \textbf{no
    original research}: an article reports what reliable sources say, it does
    not argue
\end{itemize}
\end{frame}

\begin{frame}{Using AI in this course}
\begin{itemize}
  \item \textbf{Not for writing.}  You are here to practise writing.  The
    words in your paper and your Wikipedia article must be your own.
    \begin{itemize}
      \item This includes \textbf{proofreading}: a model would rewrite your
        sentences into fluent, anonymous English.  Spell-checker yes; rewrite
        no.  For help with language, ask us.
    \end{itemize}
  \item \textbf{For research, with disclosure.}  Use it to explain a
    concept, suggest search terms, or point you towards literature.
  \item \textbf{Say what you used}, naming the tool and the version, and
    what you used it for:
    \begin{itemize}
      \item \emph{Claude Opus 5 (Anthropic) was used to suggest search
        terms.  All sources were located, read, and cited by me.}
    \end{itemize}
  \item What a model tells you is a lead, not a source: find the original,
    \textbf{read} it, and cite that.
  \item Undisclosed use counts as academic misconduct.
  \item Full policy, and how assessment works, on \website.
\end{itemize}
\end{frame}

% =====================================================================
\section{Why Academic Writing Matters}

\begin{frame}{Why It Matters}
\begin{itemize}
  \item It \textbf{helps thinking}: analysis, structure, precision
  \item It \textbf{builds knowledge}: claims + evidence + reasoning
  \item It \textbf{travels}: others can reuse, critique, extend
  \item It \textbf{signals credibility}: method, transparency, sources
  \item It \textbf{connects communities}: researchers, practitioners, public
\end{itemize}
\end{frame}

\begin{frame}{Write for your Audience}
Three questions to ask before you write:
\begin{itemize}
\item \textbf{Who are they?} \hfill terminology, the arguments they expect
\item \textbf{What do they already know?} \hfill assume that, explain the rest
\item \textbf{What do they want from you?} \hfill that depends:
\end{itemize}

\begin{center}
\begin{tabular}{@{}ll@{}}
\toprule
\textbf{Writing for} & \textbf{they want} \\
\midrule
basic research    & understanding \\
applied research  & better practice \\
literary studies  & better appreciation \\
the fine arts     & to be engaged \\
Wikipedia         & an answer, quickly \\
\bottomrule
\end{tabular}
\end{center}

\begin{itemize}
\item Different disciplines are interested in different things (not always sensibly)
\end{itemize}
\end{frame}

\begin{frame}{Define Your Purpose}
\begin{itemize}
  \item Explain, evaluate, compare, propose, synthesize?
  \item Purpose drives selection of evidence and structure
  \item One paper should have one clear purpose
    \begin{itemize}
    \item \textbf{Specific}: a claim you can support in a short paper
    \item \textbf{Contestable}: not a truism; invites argument
    \item \textbf{Roadmap}: hints at reasons/structure to come
    \end{itemize}
  \end{itemize}
\end{frame}

% =====================================================================
\section{Global Structure}

\begin{frame}{Common Structures}
\begin{itemize}
  \item Humanities: \textbf{thesis-driven essay}
  \item Social and natural sciences: \textbf{IMRaD} (+ Conclusion)
  \item Linguistics:
    \\   Intro → Background → Data/Methods → Analysis → Discussion → Conclusion
  \item Computer Science:
    \\ Intro → Related Work → Method → Data → Experiments/Results → Analysis → Conclusion
  \item Hybrids: literature review + case study; policy analysis + recommendations
\end{itemize}
\end{frame}

\begin{frame}{Start and End Well}
  \begin{itemize}
  \item The Introduction sets the stage
\begin{itemize}
  \item Context \textrightarrow\ Problem \textrightarrow\ Question/Claim \textrightarrow\ Contribution \textrightarrow\ Roadmap
  \item Avoid: history lessons without focus; claims with no stakes
\end{itemize}
\item The Conclusion drives home the argument
  \begin{itemize}
  \item Answer the question; synthesize findings
  \item State implications, limits, next steps
  \item \textbf{Avoid}: repeating the intro; new evidence in the last paragraph
  \end{itemize}
\end{itemize}
\end{frame}

\begin{frame}{A Minimal Model Outline}
\begin{itemize}
  \item \textbf{Abstract} (150--200 words)
  \item \textbf{Introduction}: context, question, thesis, roadmap
  \item \textbf{Body}: 2--3 sections building your case
  \item \textbf{Conclusion}: answer, implications, limits
  \item \textbf{References}: consistent style (we will use APA)
\end{itemize}
\end{frame}

% =====================================================================
\section{Your First Writing Task}

\begin{frame}{Compare three kinds of Writing}
\textbf{Do data centres use too much water?}
\begin{itemize}
  \item An \textbf{encyclopedia} summarises what is known
    \begin{itemize}
    \item \emph{Environmental impact of AI}, the \emph{Water usage} section
      \citep{WikiEnvAI}
    \end{itemize}
  \item An \textbf{academic paper} argues a case, hedged and cited
    \begin{itemize}
    \item \citet{Mytton2021}, \emph{Data centre water consumption} --- six
      pages, open access
    \end{itemize}
  \item A \textbf{blog post} argues a case too, but by its own rules
    \begin{itemize}
    \item \citet{Masley2025Water}, \emph{The AI water issue is fake}
    \end{itemize}
\end{itemize}
\vspace{1ex}
All three describe the same world.  They disagree about how much it matters.
\end{frame}

\begin{frame}{Write Three Paragraphs (by Sunday 11 October)}
\begin{itemize}
  \item Read all three, then write roughly 400--500 words:
    \begin{itemize}
    \item \textbf{What does each one claim?}  State it in one sentence each
    \item \textbf{Where do they differ}, and is it about the facts or about
      what the facts mean?
    \item \textbf{Which did you trust most, and why?}  Name what convinced
      you --- evidence, sources, hedging, tone
    \end{itemize}
  \item Post it on Moodle by Sunday 11 October, and bring it to class on 13 October
  \item \textbf{Your own words}: no generative AI.  Cite the three sources and anything else you read.
  \item Not graded.  We use it to see how you write; you use it to practise
    reading like a writer.  We come back to it on 13 October.
\end{itemize}
\end{frame}

\begin{frame}{What a Paragraph Should Do}
\begin{itemize}
  \item One paragraph, one point
  \item \textbf{Topic sentence first}: the reader knows the point before the detail
  \item Then support it: evidence, example, or reason
  \item End by linking to what comes next
  \item Read it aloud --- if you run out of breath, the sentence is too long
\end{itemize}
\end{frame}

% =====================================================================
\section{What's Coming}

\begin{frame}{The Rest of the Course}
\small
\begin{itemize}
  \item \textbf{29 Sep} (Pavel): Wikipedia rules and recommendations ---
    accounts, sandboxes, the five pillars in practice
  \item \textbf{13 Oct} (Francis): academic vs.\ encyclopedic writing;
    finding and evaluating sources, in any language
  \item \textbf{27 Oct} (Pavel): writing a good Wikipedia article ---
    clear, engaging, neutral
  \item \textbf{10 Nov} (Francis): making your data and writing useful and
    ethical --- FAIR, CARE, licences, open access
  \item \textbf{24 Nov} (Pavel): feedback on Wikipedia articles ---
    editing as communal knowledge building
  \item \textbf{8 Dec} (Francis): peer review and revision for publication
  \item \textbf{15 Dec}: conclusions and reflection
\end{itemize}
\vspace{1ex}
Both assignments are due in January --- dates and weights on \website.
\end{frame}

% =====================================================================
\section{(Wiki) Orientation \& Onboarding}

\begin{frame}{(Wiki) Ground Rules}
\begin{itemize}
  \item Wikipedia \textbf{Five pillars} \citep{wikipediaFivePillars}
  \item \textbf{UCoC}: Universal Code of Conduct \citep{WikimediaUCoC}
  \item \textbf{Programs \& Events Dashboard} enrollment \citep{WikiEduDashboard}
  \item Starter modules: editing basics, sourcing, plagiarism
\end{itemize}
\end{frame}

\begin{frame}{Before Next Time}
\begin{itemize}
  \item Create a Wikipedia account and enrol on the Dashboard
  \item \textbf{29 September}: Pavel's session on Wikipedia, in Czech
  \item \textbf{11 October}: post your three paragraphs on Moodle
  \item Start noticing environmental topics you might want to write about
\end{itemize}
\end{frame}

% =====================================================================
\begin{frame}[allowframebreaks]
\frametitle{References}
\printbibliography[heading=none]
\end{frame}

% =====================================================================
% Appendix: reference material, covered properly in later sessions.
\AtBeginSection[]{}
\appendix

\begin{frame}{Appendix}
These slides are here for reference.  We come back to all of them later in
the course.
\end{frame}

\begin{frame}{Choosing and Using Evidence (13 October)}
  \begin{itemize}
  \item What can be trusted as evidence?
    \begin{itemize}
    \item Peer-reviewed articles; scholarly books
    \item Policy reports; official statistics; reputable NGOs
    \item Data (quantitative/qualitative), corpora, case studies
    \end{itemize}
  \item Evaluating sources
    \begin{itemize}
    \item \textbf{Authority}: who wrote it? venue? peer review?
    \item \textbf{Recency/Relevance}: up-to-date, on-point
    \item \textbf{Method/Transparency}: can you inspect or replicate?
    \end{itemize}
  \end{itemize}
\end{frame}

\begin{frame}{Integrating Sources}
\begin{itemize}
  \item \textbf{Summarize}: key point in your words, with citation
  \item \textbf{Paraphrase}: reframe to serve your argument
  \item \textbf{Quote}: sparingly, when wording is crucial
  \item Always connect source to \textbf{your} claim
  \item Multiple sources are best
    \begin{itemize}
    \item Weave multiple sources to make a new point
    \item Compare/contrast findings, methods, assumptions
    \item Identify gaps, tensions, implications
    \end{itemize}
  \end{itemize}
\end{frame}

\begin{frame}{Multilingual Source Literacy}
\begin{itemize}
\item Often sources in other languages are ignored
  \item When English is not the richest source: local journals, government docs, linguistic data, \ldots
  \item Beware translation bias; summarize fairly
  \item Cross-check facts across languages
\end{itemize}
\end{frame}

\begin{frame}{Academic Integrity}
\begin{itemize}
\item Cite all sources (including data, images)
  \begin{itemize}
  \item Read what you cite, don't cite transitively (or acknowledge when you do)
  \item Cite in detail --- give page numbers (especially for books)
  \end{itemize}
  \item Proper paraphrase: change structure and wording, cite anyway
  \item Avoid patchwriting; keep notes disciplined
  \item Never fabricate a source --- and never trust a model's citation
    until you have found it
\end{itemize}
\end{frame}

\begin{frame}{Sustainable Knowledge (10 November)}
This is more about data than writing, but still very important.

\begin{itemize}
  \item \textbf{FAIR}: Findable, Accessible, Interoperable, Reusable \cite{FAIR2016}
  \item \textbf{CARE}: Collective benefit, Authority, Responsibility, Ethics \cite{CARE2020}
  \item \textbf{Open Science} (UNESCO Recommendation) \cite{UNESCOOpenScience2021}
  \item \textbf{Knowledge equity} (Wikimedia 2030) \cite{KnowledgeEquity2030}
\end{itemize}
\end{frame}

\begin{frame}{Writing is Revising (8 December)}
\begin{itemize}
\item Draft fast; revise for structure, then style
  \begin{itemize}
  \item \textbf{Content is more important than presentation}
  \end{itemize}
  \item Read aloud; reverse outline; get feedback early
  \item Track changes; keep a revision log
\end{itemize}
\end{frame}

\begin{frame}{Peer Review is part of the process}
  \begin{itemize}
  \item Helpful feedback
    \begin{itemize}
    \item Focus on \textbf{claims, evidence, logic}, not just grammar
    \item Ask: what is the thesis? is it supported? what is missing?
    \item Offer specific, actionable suggestions
    \item Don't be \href{https://www.youtube.com/watch?v=-VRBWLpYCPY}{reviewer three}
    \end{itemize}
  \item Responding to Feedback
    \begin{itemize}
    \item Separate criticism into categories (structure, evidence, style)
    \item Decide: change, clarify, or justify (with reasons)
    \item Write a brief \textbf{revision memo}: what changed and why
    \end{itemize}
  \end{itemize}
\end{frame}

\begin{frame}{Common Mistakes}
\small
\begin{center}
\begin{tabular}{@{}ll@{}}
\toprule
\textbf{Mistake} & \textbf{Fix} \\
\midrule
Topic too broad     & narrow to a question you can answer with the evidence you have \\
Claim without evidence & support it, and show how the evidence bears on the claim \\
Source dump         & synthesize; explain why each source matters \emph{for your thesis} \\
Structure drift     & reverse outline; re-order sections to match the argument \\
Jargon overload     & define key terms once; prefer plain English or Czech \\
\bottomrule
\end{tabular}
\end{center}
\end{frame}

\begin{frame}{Pre-Submission Checklist}
\begin{itemize}
  \item Clear thesis in the introduction
  \item Section headings reflect argumentative moves
  \item Each paragraph has a topic sentence and a purpose
  \item Claims are cited; sources are integrated (not just quoted)
  \item Figures/tables are labeled and discussed
  \item Proofread for clarity, concision, coherence
\end{itemize}
\end{frame}

\begin{frame}{What Good Work Looks Like}
\begin{itemize}
  \item \textbf{Clarity}: readers always know the claim \& why it matters
  \item \textbf{Evidence}: claims anchored to credible sources
  \item \textbf{Ethics}: accurate attribution; respectful tone; UCoC compliance
  \item \textbf{Sustainability}: work that others can find, reuse, build upon
\end{itemize}
\end{frame}

\end{document}
%%% Local Variables:
%%% coding: utf-8
%%% mode: latex
%%% TeX-PDF-mode: t
%%% TeX-engine: luatex
%%% End:
